\documentclass[a4paper,11pt]{article}
\usepackage[T2A]{fontenc}
\usepackage[utf8]{inputenc}
\usepackage[russian]{babel}
\usepackage{amsmath,amssymb,mathtools}
\usepackage{cmbright}
\renewcommand{\familydefault}{\sfdefault}
\usepackage{icomma}
\usepackage[a4paper,left=16mm,right=16mm,top=12mm,bottom=11mm]{geometry}
\usepackage{microtype}
\usepackage{multicol}
\usepackage{xcolor}
\usepackage{hyperref}
\definecolor{accent}{HTML}{15747A}
\definecolor{darkaccent}{HTML}{164C55}
\definecolor{textgray}{HTML}{283238}
\definecolor{muted}{HTML}{65747A}
\definecolor{rulegray}{HTML}{CADADC}
\color{textgray}
\hypersetup{hidelinks}
\setlength{\parindent}{0pt}
\setlength{\parskip}{0pt}
\setlength{\columnsep}{8mm}
\setlength{\abovedisplayskip}{1mm}
\setlength{\belowdisplayskip}{1mm}
\emergencystretch=2em
\newcommand{\sect}[1]{\par\vspace{2.2mm}{\normalsize\bfseries\color{darkaccent}#1}\par\vspace{0.3mm}{\color{rulegray}\rule{\linewidth}{0.45pt}}\par\vspace{0.8mm}}
\newcommand{\fml}[1]{\par\noindent{\normalsize\(\displaystyle #1\)}\par\vspace{2.8mm}}
\begin{document}
\thispagestyle{empty}

\hypersetup{pdftitle={Главные формулы. София Хоменко},pdfauthor={Лёвин Артём Александрович}}
\begin{center}
{\small\bfseries\color{accent}УЧЕБНАЯ ПАМЯТКА \quad\textbullet\quad ОСНОВНЫЕ ФОРМУЛЫ}\par
\vspace{1mm}
{\fontsize{21}{25}\selectfont\bfseries\color{darkaccent} Корни, проценты и тригонометрия\par}
\vspace{0.8mm}
{\small\color{muted}София Хоменко\quad\textbullet\quad ОГЭ · математика\quad\textbullet\quad 12.10.2026}
\end{center}
\vspace{2mm}
\noindent{\color{accent}\rule{\linewidth}{0.9pt}}
\vspace{1.3mm}
{\bfseries\color{darkaccent}Главные принципы решения}\par
\vspace{1mm}
\begin{multicols}{2}
\raggedcolumns
\small
\noindent\textbullet\quad Перед извлечением квадратного корня проверьте неотрицательность подкоренного выражения.\par\vspace{0.75mm}
\noindent\textbullet\quad При решении рационального неравенства отметьте нули и запрещённые точки.\par\vspace{0.75mm}
\columnbreak
\noindent\textbullet\quad Процент уклона — отношение высоты к горизонтальной длине, умноженное на 100\%.\par\vspace{0.75mm}
\noindent\textbullet\quad Для табличного угла опирайтесь на синус и косинус прямоугольного треугольника.\par\vspace{0.75mm}
\end{multicols}
\vspace{1.1mm}
\noindent{\color{accent}\rule{\linewidth}{0.5pt}}
\vspace{0.6mm}
\begin{multicols}{2}
\raggedcolumns
\normalsize
\sect{Корни и квадратные выражения}
\fml{\sqrt{a^2}=|a|}
\fml{\sqrt{ab}=\sqrt a\sqrt b\;(a,b\ge0)}
\fml{(a-b)(a+b)=a^2-b^2}
\sect{Геометрия и уклон}
\fml{a^2+b^2=c^2}
\fml{\text{уклон}=\frac hL\cdot100\%}
\fml{S_{\text{прям}}=ab}
\columnbreak
\sect{Тригонометрия}
\fml{\sin30^\circ=\frac12}
\fml{\cos30^\circ=\frac{\sqrt3}{2}}
\fml{\sin45^\circ=\cos45^\circ=\frac{\sqrt2}{2}}
\fml{\sin(180^\circ-\alpha)=\sin\alpha}
\sect{Проценты}
\fml{\text{часть}=\text{целое}\cdot\frac p{100}}
\fml{\text{новое}=\text{старое}\cdot\left(1-\frac p{100}\right)}
\end{multicols}
\vfill
{\color{accent}\rule{\linewidth}{0.8pt}}\par\vspace{1.5mm}
{\small\color{muted}Лёвин Артём Александрович эксклюзивно для Софии Хоменко}
\end{document}